\documentclass[11pt,draftclsnofoot,onecolumn,journal]{IEEEtran}
\usepackage{amsmath,amssymb}
\usepackage[hidelinks]{hyperref}
\usepackage{booktabs}

\newtheorem{theorem}{Theorem}
\newtheorem{lemma}{Lemma}
\newtheorem{proposition}{Proposition}
\newtheorem{remark}{Remark}
\newcommand{\ab}{\bar\alpha}

\title{An Independent Certificate for the One-Letter Bound in the
Huang--Liu--Liu Counterexample to Marton's Inner Bound}
\author{%
  [Author list TBD]\thanks{%
  Code and outputs:
  \url{https://github.com/gokhanmergen/information-theory-problems},
  directory \texttt{attempts/marton-inner-bound-tightness/code/independent-cert}.}%
\thanks{AI use: the verification code, the computations and this draft were produced by an AI
system (Claude Opus 5.5, Anthropic). The human authors are responsible for the content.}}

\begin{document}
\maketitle

\begin{abstract}
Huang, Liu and Liu (arXiv:2608.19869, not yet peer reviewed) recently claimed that Marton's
inner bound is not tight for the two-receiver broadcast channel. Their argument rests on a ternary base channel~$T$ with a
fixed-input two-letter gain of $2.83\times10^{-6}$ nats. The upper side of that gain is a
one-letter bound $UB(p^*)\le 0.0916508833913554746710$, and proving it requires certifying a
\emph{global} maximum of a non-concave information functional $G_T$ over all auxiliary
distributions. Huang--Liu--Liu did this with an interval branch-and-bound. Its stored
certificate keeps only per-face summaries and transcript hashes, so it cannot be audited
box by box.

We give an independent certificate. It uses a different design: the closed simplex with no
boundary cutoff, exact integer arithmetic for every linear form, and Arb ball arithmetic for
logarithms. It shares no numerical code with the original. It proves the bound needed by
Huang--Liu--Liu, and it also proves the sharper statement $\max G_T\in[G^\star,\,G^\star+10^{-12})$,
where $G^\star$ is the value at their reported optimizer. As a consequence, the fixed-input gap is
at least $3.8275\times10^{-6}$ nats.

We also check that the Anantharam--Gohari--Nair cardinality theorem applies to $G_T$ with explicit
parameters, and we give a self-contained proof of the determinism part. Finally, we audit the
original certificate files and reproduce all 83 of their transcript hashes bit for bit.
We found no error.
\end{abstract}

\section{Setting}

All logarithms are natural. Let $T$ be the ternary broadcast channel of~\cite[(14)--(15)]{HLL}
with marginals $T_Y,T_Z$. Let $p^*$ be the input law of~\cite[(16)]{HLL}, and let
$\alpha=\alpha_0$ and $a=(a_x)$ be as in~\cite[(17)]{HLL}. Every printed decimal is read as an
exact rational. For finite $\mathcal U,\mathcal V$ and a law $P_{UVX}$ define
\[
G(P_{UVX})=-\alpha H(Y)-\ab H(Z)+I(U;Y)+I(V;Z)-I(U;V)+\textstyle\sum_x p(x)a_x,\qquad \ab=1-\alpha,
\]
and let $B_T=\sup G$ over all finite alphabets. Huang--Liu--Liu~\cite[(23)--(24)]{HLL} show that the
Marton sum rate with common auxiliary satisfies
\[
M_T(p)\le UB(p):=\alpha H_p(Y)+\ab H_p(Z)+B_T-a\cdot p\qquad\text{for all }p\in\Delta(\mathcal X),
\]
and they assert the numerical bound $UB(p^*)\le 0.0916508833913554746710$~\cite[(25)]{HLL}. Their
optimizer~\cite[(19)]{HLL} has $X=f(U,V)$ for the map $f=\left(\begin{smallmatrix}0&1\\2&0\end{smallmatrix}\right)$
with binary $U,V$. Its value is
$G^\star:=G(19)=-1.06471340944277385029433\ldots$, and it induces exactly~$p^*$. Since
\[
UB(p^*)=I(U;Y)+I(V;Z)-I(U;V)\big|_{(19)}+(B_T-G^\star)=0.0916503833913554\ldots+(B_T-G^\star),
\]
inequality~(25) is equivalent to $B_T-G^\star\le 5.0000000000000000000189\times10^{-7}$. So the
certified quantity is the gap between the global maximum of $G$ and its value at~(19).

\section{Reduction to a finite-dimensional problem}

\begin{proposition}\label{prop:red}
$B_T$ equals the maximum of $G$ over the following domain:
\begin{itemize}
\item[(i)] $X=f(U,V)$ with $U,V$ binary, $f$ any of the $81$ maps $\{0,1\}^2\to\{0,1,2\}$, and $P_{UV}$ in the
closed $3$-simplex;
\item[(ii)] the two branches $(U=X,\ V\text{ constant})$ and $(V=X,\ U\text{ constant})$, with $p_X$ in the closed $2$-simplex.
\end{itemize}
Moreover, in (i) it suffices to take one $f$ from each of the $27$ orbits under relabelling of $U$ and of~$V$.
\end{proposition}

\begin{IEEEproof}
\emph{Cardinality.} Anantharam, Gohari and Nair~\cite[Thm.~3]{AGN} consider the maximization over
$p(x)$ of
\begin{multline*}
-(\delta_1+\lambda(\delta_0+\gamma_2))I(X;Y)-(1-\lambda)(\delta_0+\gamma_2)I(X;Z)\\
+\max_{p(u,v|x)}\big[(\delta_1+\gamma_2)I(U;Y)+\gamma_2I(V;Z)-\gamma_2I(U;V)\big]-\textstyle\sum_xd_xp(x).
\end{multline*}
They show that it is attained with $|\mathcal U|+|\mathcal V|\le|\mathcal X|+1$ and $X$ a function of $(U,V)$.
Writing $-\alpha H(Y)=-\alpha I(X;Y)-\alpha\sum_xp(x)H(T_Y(\cdot|x))$, and similarly for $Z$, shows
that $G$ is exactly this objective with
\[
\delta_0=\delta_1=0,\quad \gamma_2=1,\quad \lambda=\alpha\in(0,1),\quad d_x=-a_x+\alpha H(T_Y(\cdot|x))+\ab H(T_Z(\cdot|x)).
\]
Hence some maximizer has $|\mathcal U|+|\mathcal V|\le4$ and $X=f(U,V)$. (A maximizer exists, because
the bounds $|\mathcal U|,|\mathcal V|\le|\mathcal X|$ of~\cite{GA} make the supremum a maximum over a compact set.)

\emph{Case analysis.} This leaves $(|\mathcal U|,|\mathcal V|)\in\{(2,2),(1,3),(3,1)\}$ or smaller sizes.
\begin{itemize}
\item Smaller sizes, and $f$ undefined on zero-mass cells, are faces of the closed $(2,2)$ simplex.
$G$ is continuous there with $0\log0=0$, and every partial map extends to a full one.
\item For $(1,3)$ with $X=f(V)$, the Markov chain $V\to X\to Z$ and the fact that $X$ is a function of $V$
give $I(V;Z)=I(X;Z)$. So the value equals that of branch $V=X$. The case $(3,1)$ is symmetric.
\end{itemize}

\emph{Symmetry.} Relabelling $U$ or $V$ leaves every term of $G$ invariant. The resulting group of
order~$4$ has $(81+9+9+9)/4=27$ orbits on the maps.
\end{IEEEproof}

\begin{remark}
The determinism part has a two-line proof that does not depend on~\cite{AGN}. Fix $P_{UV}$. The
maps $P_{X|UV}\mapsto P_{Y|U},P_{Z|V},p_Y,p_Z,p_X$ are linear, and each term of $G$ is convex in
$P_{X|UV}$:
\begin{itemize}
\item $I(U;Y)$ is convex in $P_{Y|U}$ for fixed $P_U$, and likewise $I(V;Z)$;
\item $-H(Y)$ and $-H(Z)$ are convex;
\item $a\cdot p_X$ is linear and $I(U;V)$ is fixed.
\end{itemize}
So the maximum over $P_{X|UV}$ is attained at a vertex, which is a deterministic map. The
cardinality bound $|\mathcal U|+|\mathcal V|\le4$ is the only input taken from~\cite{AGN}. We read the
proof of~\cite[Prop.~3, Lemma~1]{AGN} but did not re-derive its appendix lemma on linearity of
$H(V|Z)$ under multiplicative perturbations.
\end{remark}

\section{The certificate}

For a map $f$ and the cell vector $p\in\Delta_3$,
\[
G(p)=\ab H(Y)+\alpha H(Z)+H(UV)-H(UY)-H(VZ)+a\cdot p_X .
\]
This is a sum of $w\,h(c\cdot p)$ terms with $h(t)=-t\log t$ (concave), $w>0$ and $c\ge0$; of
$w\,g(d\cdot p)$ terms with $g(t)=t\log t$ (convex); and of a linear term. The two branches have the
same form on $\Delta_2$.

\begin{lemma}[box bound]\label{lem:box}
Let $K=\{lo\le p\le hi,\ \sum_ip_i=1\}$, and let $[t_1,t_2]$ be the exact range of $c\cdot p$ on $K$.
\begin{itemize}
\item For every $m>0$ and every $t\ge0$, $h(t)\le m-(1+\log m)\,t$.
\item For every $A\ge g(t_1)$ and $B\ge g(t_2)$, the chord through $(t_1,A)$ and $(t_2,B)$ dominates $g$ on $[t_1,t_2]$.
\end{itemize}
Replacing each coefficient of the resulting affine function of $p$ by an upper bound (legal since
$p\ge0$) and maximizing over $K$ by the greedy rule gives an upper bound on $\sup_KG$.
\end{lemma}

The tangent bound holds at $t=0$, so no cutoff away from the boundary of the simplex is needed.
This differs from~\cite{HLL}, who certify a $\delta$-interior for $\delta=2^{-30}$ and handle the
boundary by recursion over faces.

\emph{Arithmetic.} Box corners are integers over $2^{160}$ and channel coefficients are integers
over $10^{12}$. Every range $[t_1,t_2]$ and every greedy step is therefore computed exactly with
integers. Logarithms, and products and quotients involving them, are evaluated in Arb~\cite{Arb}
via python-flint at 128 bits. Every quantity entering the bound is the upper endpoint of a ball,
and the test ``bound $<$ target'' is a certain ball comparison.

\emph{Search.} Depth-first branch-and-bound. A box is discarded if its bound is below the target.
Otherwise its bounds-consistent hull is bisected along its longest side. Coverage holds by
construction.

\emph{Targets.} We take $T_5$ to be a rational on the $10^{-30}$ grid (the largest one below
the lower endpoint of an Arb enclosure, hence conservative) such that
\[\alpha H_{p^*}(Y)+\ab H_{p^*}(Z)+T_5-a\cdot p^*\le 0.0916508833913554746710,\]
computed with outward rounding. This gives $T_5-G^\star=5.0000000000000000189\times10^{-7}$. The sharper target is
$T_{12}=-10647134094417738503/10^{19}$, for which $T_{12}-G^\star=9.99999994\times10^{-13}$.

\begin{theorem}\label{thm:main}
$G<T_{12}$ on the whole domain of Proposition~\ref{prop:red}. Hence
\[
G^\star\le B_T<G^\star+10^{-12},
\]
and in particular $B_T<T_5$, so~\cite[(25)]{HLL} holds. Moreover $UB(p^*)\le 0.09165038339235547467$.
\end{theorem}

The proof is computer-assisted. Table~\ref{tab:runs} lists the run data; both runs used two
processes on an Apple-silicon laptop.

\begin{table}[h]
\centering
\caption{Branch-and-bound runs (27 map orbits plus 2 branches).}\label{tab:runs}
\begin{tabular}{lrrrr}
\toprule
target & boxes & boxes for HLL's map & max depth & wall time\\
\midrule
$T_5$ & 2{,}531{,}505 & 459{,}967 & 52 & 406 s\\
$T_{12}$ & 3{,}110{,}939 & 1{,}033{,}045 & 86 & 319 s\\
\bottomrule
\end{tabular}
\end{table}

\emph{Controls.} With target $G^\star-10^{-9}$ the search must fail, and it does: it stops at an
unsplittable box within $2\times10^{-4}$ of~(19). A randomized test compares Lemma~\ref{lem:box}'s
bound with $G$ evaluated directly from its mutual-information form at 50 digits. It covers 400
boxes with widths from $2^{-1}$ to $2^{-45}$, 30\% of them touching the boundary, and finds no violation.

\section{Consequence for the fixed-input gap}

Take Huang--Liu--Liu's explicit two-letter law~\cite[App.~A]{HLL}, with $W=(Q,W_c)$ and $P(Q=0)=\mu$ as
printed. Each of its two components has both coordinate marginals exactly equal to $p^*$, which
is~\cite[(26)]{HLL}. Since $I(Q,W_c;Y^2)\ge I(W_c;Y^2\mid Q)$ and the conditional terms average
exactly over $Q$, Arb evaluation gives
\[
M_{T^{\otimes2}}(P^*)\ \ge\ V_{\rm lo}=0.18330459432065522490176\ldots\ .
\]
By Theorem~\ref{thm:main},
\[
V_{\rm lo}-2\,UB(p^*)\ \ge\ 3.8275\times10^{-6}\ \text{nats}.
\]
With the bound~(25) alone the gap is $2.8275\times10^{-6}$, as printed in~\cite{HLL}. This is the
fixed-input input to~\cite[Thm.~4]{HLL}, which then yields an unconstrained finite channel with
$M_{\tilde T^{\otimes2}}>2M_{\tilde T}$. We did not re-check~\cite[Thms.~2,~4]{HLL}, the gradient-shaping
and constraint-removal steps, beyond reading them.

\section{Audit of the original certificate}

We read the original certificate files (seven JSON files) as untrusted data and checked them
with our own Arb code. Every check passed:
\begin{itemize}
\item the instance equals the paper's;
\item the stored target $B$ satisfies $B-G^\star=5\times10^{-7}$;
\item the stored boundary modulus dominates its formula;
\item all $81$ maps and both branches are present, each with all $15$ (respectively $7$) faces;
\item face targets equal $B-(4-|S|)\omega$ exactly;
\item every recorded pruned upper bound is below its target, and no boxes remain;
\item the vertex values agree with ours to $10^{-40}$;
\item the box total is $4{,}729{,}704$.
\end{itemize}

The individual boxes are not stored, so the bounds cannot be re-derived from the files. Instead,
we re-ran the original engine under different gmpy2 and MPFR versions, and all $83$ transcript
digests matched. We also prove a boundary lemma that justifies their face recursion with a
sharper modulus. Let $p=(1-r)q+re_i$, where $q$ puts no mass on cell~$i$. With $B$ the indicator of
cell~$i$,
\begin{align*}
G(p)&=\textstyle\sum_bP(b)G(p_b)+\ab I(B;Y)+\alpha I(B;Z)+H(B)-I(B;UY)-I(B;VZ)\\
&\le(1-r)G(q)+rG(e_i)+H_2(r).
\end{align*}

\section{What is and is not established}

Theorem~\ref{thm:main} rests on four things:
\begin{itemize}
\item the cited theorem~\cite[Thm.~3]{AGN};
\item Arb's guaranteed enclosures;
\item exact integer arithmetic;
\item the correctness of about $350$ lines of our code, whose soundness argument is
Lemma~\ref{lem:box}.
\end{itemize}
The code has not yet been reviewed by a second party. We did not verify~\cite[Thms.~2,~4,
Prop.~3]{HLL} or the explicit lifted instance~\cite[(30)--(43)]{HLL}. An earlier non-rigorous
recomputation in our repository reproduces~(40)--(43).

\begin{thebibliography}{9}
\bibitem{HLL} M.~Huang, Y.~Liu, and Y.~Liu, ``Sub-optimality of Marton's inner bound for the
two-receiver broadcast channel,'' arXiv:2608.19869v1, Aug.~2026; code at
\url{https://github.com/yanxiaoliu-mike/Suboptimality_Marton}.
\bibitem{AGN} V.~Anantharam, A.~Gohari, and C.~Nair, ``On the evaluation of Marton's inner bound
for two-receiver broadcast channels,'' \emph{IEEE Trans. Inf. Theory}, vol.~65, no.~3,
pp.~1361--1371, 2019.
\bibitem{GA} A.~Gohari and V.~Anantharam, ``Evaluation of Marton's inner bound for the general
broadcast channel,'' \emph{IEEE Trans. Inf. Theory}, vol.~58, no.~2, pp.~608--619, 2012.
\bibitem{Arb} F.~Johansson, ``Arb: efficient arbitrary-precision midpoint-radius interval
arithmetic,'' \emph{IEEE Trans. Comput.}, vol.~66, no.~8, pp.~1281--1292, 2017.
\end{thebibliography}

\end{document}
